\chapter{测试用例设计}

为了验证 LingoBridge 中文学习平台在各类复杂业务场景下的行为正确性与系统稳定性，测试团队严格遵循 GB/T 9385—2008《计算机软件需求规格说明规范》\cite{gb9385} 与 IEEE 1016-2009《软件设计描述标准》\cite{ieee1016} 中关于测试可追溯性与功能覆盖的工程要求，设计了一套覆盖系统全生命周期核心功能模块的测试用例矩阵。测试用例的设计综合应用了黑盒测试方法中的等价类划分法、边界值分析法以及端到端业务场景法，特别针对高频交互和强数据依赖的模块进行了深度覆盖。测试设计的核心目标是确保在各种正常及异常输入下，系统均能给出符合设计预期的正确响应，并将异常阻断在用户交互界面之外。

以下测试用例矩阵详尽列出了针对用户身份认证与登录、课程目录大纲展示、交互式富媒体课件播放与进度持久化、课后语音作业提交与 AI 智能打分、高并发直播互动信令交互以及管理后台报表生成等关键功能组件的测试设计，以此作为自动化测试与手工测试实施的具体执行指南。

\begin{longtable}{@{}p{2.2cm} p{1.6cm} p{4.6cm} p{4.6cm} p{1.2cm}@{}}
\caption{LingoBridge 核心功能模块测试用例设计矩阵} \label{tab:test-cases} \\
\toprule
用例编号 & 测试模块 & 测试步骤与输入 & 预期结果 & 优先级 \\
\midrule
\endfirsthead
\multicolumn{5}{r}{续表 \ref{tab:test-cases}} \\
\toprule
用例编号 & 测试模块 & 测试步骤与输入 & 预期结果 & 优先级 \\
\midrule
\endhead
\bottomrule
\endfoot
\bottomrule
\endlastfoot
TC\_Lingo\_001 & 登录模块 & 用户在前端登录界面中输入合法的注册用户名与正确密码，随后点击登录按钮触发身份认证请求。 & 系统成功完成身份校验并返回代表用户登录状态的 JSON Web Token，前端自动跳转至学员个人主页并渲染个性化推荐课程。 & 高 \\
\midrule
TC\_Lingo\_002 & 登录模块 & 用户输入不存在的用户名或者错误的密码尝试登录，随后点击登录按钮触发身份认证流程。 & 系统拒绝该认证请求并返回特定的错误码，前端界面友好展示账户或密码错误的提示信息，且不允许进行后续任何敏感 API 的调用。 & 高 \\
\midrule
TC\_Lingo\_003 & 课程模块 & 登录成功的用户点击进入课程中心，选择特定的中文学习课程，并点击展开课程知识树大纲。 & 系统能够秒级加载并渲染出完整的课程知识图谱与目录结构，准确显示每个子课件的已完成或未开始的学习状态。 & 中 \\
\midrule
TC\_Lingo\_004 & 课件播放 & 用户打开包含富媒体视频的交互式课件，在播放五分钟后手动关闭浏览器窗口以模拟中途异常退出。 & 前端播放器通过事件监听自动捕获播放断点，并将当前的播放进度时间戳通过心跳 API 实时且安全地持久化到后端数据库中。 & 高 \\
\midrule
TC\_Lingo\_005 & 在线作业 & 用户进入课后作业界面并点击开启语音答题，调用浏览器底层麦克风录制一段标准的中文普通话朗读音频后点击提交。 & 前端音频组件实时编码上传音频流，智能语音打分微服务成功对其进行多维度声学特征分析，并在三秒内返回包含发音分值与纠错建议的详细评测报告。 & 高 \\
\midrule
TC\_Lingo\_006 & 直播互动 & 在高并发直播课程开始时，多名在线学生同时进入虚拟直播交互室，发送文字弹幕并尝试向讲师发起连麦申请。 & 直播信令服务器成功处理并发连接，文字弹幕在所有在线客户端实现无延迟渲染，音视频推拉流在 WebRTC 架构下保持极低时延与画面同步。 & 中 \\
\midrule
TC\_Lingo\_007 & 管理后台 & 管理员使用高级权限账号登录管理后台，点击导出本月全体学员的中文学习进度、答题历史与智能语音评测统计报表。 & 后端报表服务在后台高效执行多表关联查询与数据聚合，成功生成并下载格式规整且内容准确的 Excel 电子表格文件。 & 低 \\
\bottomrule
\end{longtable}
